\section{Proofs for Analysis of $\calF_{\LMKPM}$}

\subsection{Proof of Correctness of $\dictInfer$}
\label{sec:dictinfer:correctness}


\begin{lemma}
	\label{lem:dictinfer_MR}
	Let $X$ and $Y$ be two sets of records and let $(\hXone, \hYone)$ be the
        hidden shuffled records received by $\tP_1$.  Moreover, let
        $f_1: \Ispace{} \to \Uspace$ be the hiding function sampled in
        $\calF_\LMKPM$ to produce $\hXone$ and $\hYone$.  If $D$ is a dictionary
        mapping $f_1(v)$ to $v$ for all $v \in \Ispace[X]$, then
        $\dictInfer(D, \hYone)$ (Figure~\ref{fig:dict_infer}) produces an
        $X$-reconstruction of $Y$.
\end{lemma}
\begin{proof}
	Let $R$ be the multiset produced by \dictInfer.
	Note that for each $\hby \in \hYone$ exactly one reconstructed record
        $\br$ is added to $R$. Let $d: R \to \hYone$ denote the function mapping
        these reconstructed records to the corresponding $\hby$. Clearly, $d$ is
        a bijection.

	We now define an injective function $\varphi: R \to Y$ and prove that it
        satisfies the three conditions for $X$-Reconstruction (Definition
        \ref{def:max_recons}). For any record $\br \in R$ let
        $(d_1, \dots, d_k)=d(\br) \in \hYone$ and let
	$$\varphi(\br) := (f_1^{-1}(d_1), \dots, f_1^{-1}(d_k))$$
	
	Note that $f_1^{-1}(d_j)$ exists for all $\br$ and $j$, since $d(\br) \in \hYone$ by definition.
	Moreover, since $f_1$ is injective, $f_1^{-1}$ is injective for elements in $f_1(\Ispace)$.
	$\varphi$ is injective since $d$ and $f_1^{-1}$ are injective.
	We show the three properties from Definition~\ref{def:max_recons} separately:
	\begin{enumerate}
		\item Let $\br \in R$, $k := |\br|$ and $i \leq |\br|$. 
		Moreover, let $(d_1, \dots, d_{k'}) = d(\br)$ for some $k' \geq k$.
		By definition of $d$ and by construction of $\br$ (line~\ref{lin:substitute_add_id}), 
		there exists some $j^*\in[k']$ such that $r[i] = D[d_{j^*}]$.
		Suppose for a contradiction that $\br[i] \not \in \varphi(\br)$, that is,
		we have $\br[i] \neq f_1^{-1}(d_j)$ for all $j \in [k']$.
		But since for all $j \in [k']$ with $d_j \in D$ we have $f_1^{-1}(d_j) = D[d_j]$,
		we also have that $\br[i] \neq D[d_j]$, i.e., there is no $j$ such that $\br[i] = D[d_j]$. 
		This is a direct contradiction to $\br[i] = D[d_{j^*}]$.
		\item Let $\br\in R$ and $(d_1, \dots, d_{k'}) = d(\br)$ for some $k' \in \N$.
		Let $v \in \varphi(\br)$, 
		i.e., there is some $i \in [k']$ such that $v = f_1^{-1}(d_i)$.
		If $v \in \Ispace[X]$, then by assumption we have $f_1(v) \in D$ and $D[f_1(v)] = v$.
		Since $v = f_1^{-1}(d_i)$, we have $f_1(v) = f_1(f_1^{-1}(d_i)) = d_i$ 
		and thus $d_i \in D$. Therefore, the condition on line~\ref{lin:substitute_membership_check} is met and
		$D[d_i] = D[f_1(v)] = v$ is added to $\br$.
		\item Holds trivially, since we have $\varphi(R) = Y$ by the definition of $d$. 
	\end{enumerate}
The lemma thus follows.
\end{proof}

\subsection{Pseudocode of $\recenumattack$}
\label{sec:recenumattack:pseudocode}
The pseudocode of $\recenumattack$ can be found in Figure~\ref{fig:rec_enum_attack}.


\subsection{Proof of Correctness of $\recenumattack$}
\label{sec:recenumattack:correctness}
\begin{proof}
In order to prove the theorem we prove the following proposition. 
The theorem follows directly from Proposition~\ref{thm:recenum_D} and Lemma~\ref{lem:dictinfer_MR}.
\begin{proposition}\label{thm:recenum_D}
	Let $T$ and $Y$ be two sets of records and let
	$D$ be the dictionary constructed during an execution of\newline$\recenumattack^{\calF_\LMKPM(\cdot, Y)}(T)$ (Figure~\ref{fig:rec_enum_attack}).
	Moreover, let $f_1$ be the hiding function sampled during the evaluation of $\calF_\LMKPM$ in line~\ref{lin:recenum_eval} in Figure~\ref{fig:rec_enum_attack}. 
	We have $D[f_1(v)] = v$ for all $v \in \Ispace[T]$.
\end{proposition}
\begin{proof} \textbf{of Proposition~\ref{thm:recenum_D}} 
	Let $X = \{\bx_0, \dots, \bx_{n-1}\}$, where $\bx_i$ denotes the altered target record constructed in 
	line~\ref{lin:recenum_extend} for $i \in \{0, \dots, n - 1\}$.
	We show that $D[f_1(v)] = v$ for all $v \in \Ispace[X]$.
	Since $\Ispace[T] \subseteq \Ispace[X]$, this implies the claim.
	
	We first show that $u_0 = f_1(v_0)$ for the value $u_0$ defined on line~\ref{lin:recenum_e0}.
	Note that the transformation of an index's binary representation to a sequence of $v_0$ and $v_1$
	in lines \ref{lin:idx_to_seq_bin_encode_start}~-~\ref{lin:idx_to_seq_bin_encode_end} of \IdxToSeq{}
	could produce two encodings that only contain one value, namely, $(v_0)^\ell$ and $(v_1)^\ell$
	for $i = 0$ and $i = 2^\ell - 1$ respectively.
	However, the latter is prevented by the condition on line~\ref{lin:seq_to_idx_tiebreak} and therefore
	\IdxToSeq{} only outputs one sequence where all values are equal.
	Hence, there is only one record in $X$ whose first $\ell$ values all coincide,
	namely, $\bx_0$ with encoding $(v_0)^\ell$.
	Since $f_1$ is injective and the values in records of $X$ are not shuffled, 
	the latter also holds for $\hXone$.
	Therefore, $\hbx_0 = (f_1(\bx_0[1]), \dots, f_1(\bx_0[\lambda + \ell]))$ (line~\ref{lin:recenum_find0}).
	Since $\bx_0[1] = v_0$, we have $u_0 = \hbx_0[1] = f_1(\bx_0[1]) = f_1(v_0)$ (line~\ref{lin:recenum_e0}).
	
	We now show that the attack correctly recovers indices from $\hXone$. 
	We have already established that $\hbx_0$ is correctly matched with $\bx_0$,
	as it's encoding is the only one containing only $v_0$. 
	Any other encoding $\bs$ produced by \IdxToSeq{}
	either (a) is $(v_1)^{\ell - 1} \| v_2$ or (b) only
	contains the values $v_0$ and $v_1$.
	\begin{description}
	\item[Case (a):] This case is only produced on input $i = 2^\ell - 1$ and the resulting vector of hidden values in the 	
		 leakage is $(f_1(v_1))^{\ell - 1} \| f_1(v_2)$. 
		 This satisfies the condition on line~\ref{lin:seq_to_idx_tiebreak} of \SeqToIdx{}, 
		 which therefore yields the correct index $2^\ell - 1$.
	\item[Case (b):] In this case, the vector of hidden values corresponding to $\bs$ in the leakage, which we denote
		$\hbs:=(f_1(\bs[1]),\allowbreak\dots,f_1(\bs[\ell]))$, does not satisfy the condition on line~\ref{lin:seq_to_idx_tiebreak}, 
		since the only encoding that would do so is $(v_1)^{\ell}$. 
		But this is never output by \IdxToSeq{}. 
		Therefore, $\hbs$ is interpreted as a binary encoding, where every occurrence of $u_0$ 
		is interpreted as $0$ and every other hidden value as $1$.
		The correctness of this recovery therefore follows from the fact that $u_0 = f_1(v_0)$,
		which we established above.
	\end{description}
	
	By the definition of $\calF_\LMKPM$, we have $\hbx = (f_1(\bx_i[1]), \dots, f_1(\bx_i[\lambda+\ell]))$ for all $i \in \{0, \dots, n-1\}$.
	Since $\Ispace[X] := \bigcup_{\substack{l}{\bx\in X\\v \in \bx}} v$, 
	the claim holds by the construction of $D$ in lines~\ref{lin:recenum_dict_start}~-~\ref{lin:recenum_dict_end}.
\end{proof}
	
% \begin{proof}
	The theorem follows directly from Proposition~\ref{thm:recenum_D} and Lemma~\ref{lem:dictinfer_MR}.
\end{proof}

\subsection{Pseudocode of $\snakeattack$}
\label{sec:snakeattack:pseudocode}

The pseudocode of $\recenumattack$ can be found in Figure~\ref{fig:snakeattack}.



\subsection{Proof of Correctness of $\snakeattack$}
\label{sec:snakeattack:correctness}


\begin{proof}
	To prove the theorem, we prove the following proposition.
	The theorem follows directly from Proposition~\ref{thm:snake_D} and Lemma~\ref{lem:dictinfer_MR}.
\begin{proposition}
	\label{thm:snake_D}
	Let $T$ and $Y$ be two sets of records and let $D$ be the substitution dictionary
	constructed during an invocation of $\snakeattack^{\calF_\LMKPM(\cdot, Y)}(T)$.
	Moreover, let $f_1$ be the hiding function sampled during the invocation of $\calF_\LMKPM$
	in line~\ref{lin:snake_evaluation}.
	Then, we have $D[f_1(v)] = v$ for all $v \in \Ispace[T]$.
\end{proposition}
\begin{proof} \textbf{of Proposition~\ref{thm:snake_D}}
	Let $X := \{\bx_1, \dots, \bx_n\}$ where $\bx_i$ is produced by \snakeencodeUnique{} in line~\ref{lin:snake_encode}
	and let $k' := |\bx_i|$ for any $i\in [n]$\footnote{We assume that all records have the same length $\lambda$.
	Thus, $k' \in \{\lambda + 1, \lambda+2\}$}.
	We show that $D[f_1(v)] = v$ for all $v \in \Ispace[X]$.
	Since \snakeattack{} potentially adds fresh values to records of $T$, but does not remove any, 
	we have $\Ispace[T]\subseteq\Ispace[X]$, which implies the claim.
	
	Similar to the proof of Proposition~\ref{thm:recenum_D}, we show that
	$\hbx_i[j] = f_1(\bx_i[j])$ for all $i\in[n]$ and all $j \in [k']$.
	The result then follows from the construction of $D$ in lines~\ref{lin:snake_D_start}~-~\ref{lin:snake_D_end}. 
	In particular, we will show that the linked list we create in \IdxToSeq{} is preserved in the leakage, 
	and that this linked list allows us to uniquely identify each record. 
	We show this by induction on $i$. 
	Let $j^*$ denote the index of a column whose values are all unique (note that if such a column does not exist, it is created on line~\ref{lin:distinctIDs_j*}).

	For $i = 1$, observe that the value $t_1[j^*]$ is not appended to any record in \snakeencodeUnique{}; 
	in other words, it is the only item in the linked list that doesn't have anything pointing to it.
	In contrast, for all $i' > 1$, the value $t_{i'}[j^*]$ is appended to $t_{i'+1}$.
	$\hbx_1$ is therefore the only record in $X$ whose $j^*$-th value is not contained in the last column, 
	i.e., $\hbx_1[j^*] \not \in \{\hbx_{i'}[k'] \setdsc i'\in[n]\}$.
	By the definition of $\calF_\LMKPM$, this property is preserved in the leakage $\hXone$.
	That is, the corresponding record in $h = (f_1(\hbx_1[1]), \dots, f_1(\hbx_1[k']))$
	is the only record in $\hXone$ such that $h[j^*]\not\in\{h_{i'}[k'] \setdsc i' \in [n]\}$.
	In line~\ref{lin:snake_recover_c} in \snakerecoverUnique{} we pick $h_\ind := h$.
	Thus, for all $j \in [k']$ we have that $h_\ind[j] = f_1(\bx_1[j])$
	and, therefore  $\hbx_1[j] = h_\ind[j] = f_C(\bx_1[j])$, where the first equality follows from line~\ref{lin:snake_recover_assign_t_star}.

	Now assume that for some $i < n$ we have $\hbx_i[j] = f_1(\bx_i[j])$ for all $j \in [k']$. 
	We show that there is a unique record which $\bx_i[k']$ is pointing to (that has not been traversed by the previous items in the linked list), 
	and that, once again this property is preserved in the leakage allowing for unique identification of the record.
	By construction in \snakeencodeUnique{}, we have $\bx_{i+1}[j^*] = \bx_i[k']$.
	Since all values in the $j^*$-th column are distinct, 
	$\bx_{i+1}$ is the only record in $X$ that satisfies this condition.
	As before, since $f_1$ is injective, the corresponding record in the leakage 
	$h = (f_1(\bx_{i+1}[1]), \dots, f_1(\bx_{i+1}[k']))$
	is the only record that satisfies $h[j^*] = \hbx_i[k']$.
	Moreover, all hidden values in the $j^*$-th column of $\hXone$	are distinct. % $t'_{i+1}$ is the only record in $\setT'$
	Thus, the dictionary $\loc$ in \SeqToIdx{} contains exactly the $n$ keys $\{f_1(\bx_{i'}[j^*]) \setdsc i' \in [n]\} = \{f_1(\hbx_{i'}[j^*]) \setdsc i' \in [n]\}$,
	which are mapped as $\loc[f_1(\bx_{i'}[j^*])] = i'$.
	In line~\ref{lin:snake_recover_assign_t_star}, we pick $\hbx_{i+1} := h_\ind$ for $\ind = \loc[\hbx_i[k']]$, 
	implying $\hbx_{i+1}[j^*] = \hbx_i[k']$ and, thus, $\hbx_{i+1} = h$.
	We therefore have $\hbx_{i+1}[j] = h[j] = f_1(\bx_{i+1}[j])$ for all $j \in [k']$.
\end{proof}

	The statement follows from Proposition~\ref{thm:snake_D} and Lemma~\ref{lem:dictinfer_MR}.
\end{proof}
